\begin{thebibliography}{99}

\bibitem[Aghion et al.(2015)]{AghionEtAl2015}
Aghion, P., Cai, J., Dewatripont, M., Du, L., Harrison, A., and Legros, P. (2015).
Industrial policy and competition.
\emph{American Economic Journal: Macroeconomics}, 7(4), 1--32.

\bibitem[Athey et al.(2013)]{AtheyCoeyLevin2013}
Athey, S., Coey, D., and Levin, J. (2013).
Set-asides and subsidies in auctions.
\emph{American Economic Journal: Microeconomics}, 5(1), 1--27.

\bibitem[Chiappinelli et al.(2025)]{ChiappinelliGiuffridaSpagnolo2025}
Chiappinelli, O., Giuffrida, L. M., and Spagnolo, G. (2025).
Public procurement as an innovation policy: Where do we stand?
\emph{International Journal of Industrial Organization}, 100, 103157.

\bibitem[Eckersley et al.(2023)]{EckersleyEtAl2023}
Eckersley, P., Flynn, A., Lakoma, K., and Ferry, L. (2023).
Public procurement as a policy tool: The territorial dimension.
\emph{Regional Studies}, 57(10), 2087--2101.

\bibitem[Justman et al.(2005)]{JustmanThisseVanYpersele2005}
Justman, M., Thisse, J.-F., and van Ypersele, T. (2005).
Fiscal competition and regional differentiation.
\emph{Regional Science and Urban Economics}, 35(6), 848--861.

\bibitem[Li(2005)]{Li2005}
Li, X. (2005).
A note on expected rent in auction theory.
\emph{Operations Research Letters}, 33(5), 531--534.

\bibitem[Li and Zheng(2009)]{LiZheng2009}
Li, T. and Zheng, X. (2009).
Entry and competition effects in first-price auctions: Theory and evidence from procurement auctions.
\emph{Review of Economic Studies}, 76(4), 1397--1429.

\bibitem[Lin and Li(2026)]{LinLi2026}
Lin, C. and Li, Y. (2026).
A unified national market and local industrial policy competition.
\emph{Frontiers of Economics in China}, 21(1), 1--32.

\bibitem[Marion(2007)]{Marion2007}
Marion, J. (2007).
Are bid preferences benign? The effect of small business subsidies in highway procurement auctions.
\emph{Journal of Public Economics}, 91(7--8), 1591--1624.

\bibitem[Motta and Polo(2026)]{MottaPolo2026}
Motta, M. and Polo, M. (2026).
Supply chain disruption and precautionary industrial policy.
\emph{International Journal of Industrial Organization}, 104, 103239.

\bibitem[Macatangay(2016)]{Macatangay2016}
Macatangay, R. E. M. (2016).
Optimal local content requirement policies for extractive industries.
\emph{Resources Policy}, 50, 244--252.

\bibitem[OECD(2026)]{OECD2026}
OECD (2026).
\emph{Public Procurement, Trade and Industrial Policies: Supporting National Priorities}.
OECD Publishing, Paris.

\bibitem[Saito and Mitsuta(2012)]{SaitoMitsuta2012}
Saito, T. and Mitsuta, N. (2012).
Analyzing local government procurements in Japan: Review of regional requirements.
\emph{Studies in Regional Science}, 42(2), 363--376.

\bibitem[Slattery(2025)]{Slattery2025}
Slattery, C. (2025).
Bidding for firms: Subsidy competition in the United States.
\emph{Journal of Political Economy}, 133(8), 2563--2614.

\bibitem[Szech(2011)]{Szech2011}
Szech, N. (2011).
Optimal advertising of auctions.
\emph{Journal of Economic Theory}, 146(6), 2596--2607.

\bibitem[Vagstad(1995)]{Vagstad1995}
Vagstad, S. (1995).
Promoting fair competition in public procurement.
\emph{Journal of Public Economics}, 58(2), 283--307.

\bibitem[Van Biesebroeck and Zhang(2026)]{VanBiesebroeckZhang2026}
Van Biesebroeck, J. and Zhang, H. (2026).
Lessons on industrial policy from industrial organization: Introduction to the special issue.
\emph{International Journal of Industrial Organization}, 104, 103243.

\bibitem[Xu and Li(2019)]{XuLi2019}
Xu, M. and Li, D. Z. (2019).
Equilibrium competition, social welfare and corruption in procurement auctions.
\emph{Social Choice and Welfare}, 53, 443--465.

\bibitem[Anonymous(2026)]{Companion2026}
Anonymous (2026).
Industrial policy composition and regional value chains.
Companion manuscript.

\end{thebibliography}
